\chapter{FUTURAS LÍNEAS DE TRABAJO}

El alcance del TFM está acotado al diseño, la implementación y la validación de un \textit{pipeline} reproducible de comparación de métricas de similitud espectral. El trabajo deja abiertas siete líneas de continuación:

\begin{itemize}
    \item \textbf{Cobertura de los polímeros ausentes en la biblioteca de referencia.} La ejecución del experimento reveló que BioPBS y ECOVIO no figuran en el catálogo de la \emph{Hagelskjær Microplastic Solution Library~1.1}, de modo que identificarlos es imposible por construcción. Combinar esa biblioteca con otras públicas, o incorporar espectros de referencia propios de ambos polímeros, es lo único que permitiría cubrir los cinco materiales medidos.
    \item \textbf{Explotación de la región \textit{fingerprint}.} El estudio se ha restringido a la región C--H ($\sim$2800--3100~cm$^{-1}$). Ya se dispone de espectros de los mismos cinco polímeros en la región \textit{fingerprint} ($\sim$100--1800~cm$^{-1}$), aportados por el director del proyecto. Parametrizando la ventana de recorte del \textit{pipeline} se podría repetir el análisis en esa región y en la combinación de ambas; la hipótesis es que la mayor riqueza de bandas discrimine mejor entre polímeros.
    \item \textbf{Contraste entre la referencia correcta y las incorrectas.} Hasta ahora la robustez se mide de forma binaria: el mejor resultado acierta o no acierta. Medir además la distancia entre la similitud con la referencia correcta y la que obtienen las referencias equivocadas daría una comparación más fina entre métricas, y convertiría el hallazgo de cobertura en evidencia cuantitativa. Los datos ya generados bastan para hacerlo, sin repetir el experimento.
    \item \textbf{Ampliación a bases de datos de referencia de mayor tamaño.} El motor de \textit{matching} distribuido con Apache Spark escala el número de espectros de referencia sin rediseñar el \textit{pipeline}. El paso obvio es repetir el análisis con bibliotecas públicas más amplias, como Open Specy o RRUFF, para las que \texttt{io.py} ya tiene carga multiformato.
    \item \textbf{Despliegue en infraestructura distribuida real.} La paralelización con PySpark se ha validado en un entorno de un único nodo. Extender la validación a un clúster multi-nodo, \textit{on-premise} o en \textit{cloud}, daría la curva de escalado real frente al número de \textit{workers}.
    \item \textbf{Extensión a más polímeros y a espectroscopía FTIR.} Con más familias de polímeros y con espectros FTIR se vería si el \textit{ranking} de métricas observado aquí se mantiene entre técnicas espectroscópicas distintas.
    \item \textbf{Incorporación de nuevas métricas o de métricas aprendidas.} El módulo \texttt{metrics.py} es un conjunto de funciones intercambiables: añadir métricas clásicas nuevas, o métricas aprendidas mediante \textit{embeddings} entrenados, es directo, y todas quedarían comparadas bajo el mismo marco experimental y estadístico ya validado.
\end{itemize}

Con más polímeros y más bibliotecas, si las diferencias entre métricas llegan a confirmarse, los resultados dan para una publicación conjunta: hoy no existe en la literatura una guía empírica de qué métrica de similitud conviene según la calidad esperada del espectro.
